\hypertarget{chapter-07-page-155}{}
\subsection{奇异积分的带权不等式}
\label{sec:ch7-weighted-singular-integrals}

本节证明加权范数不等式. 回忆定义~\ref{def:ch5-calderon-zygmund-operator}: Calderón--Zygmund 算子是在 $L^2$ 上有界, 并且对 $f\in\mathcal S$ 可表示为
\[
 Tf(x)=\int_{\mathbb R^n}K(x,y)f(y)\,dy,
 \qquad x\notin\operatorname{supp}(f),
\]
的算子, 其中 $K$ 是标准核, 即满足 \eqref{eq:ch5-standard-kernel-size}, \eqref{eq:ch5-standard-kernel-y-regularity} 和 \eqref{eq:ch5-standard-kernel-x-regularity} 的核.

先证明两个引理.

\begin{Lemma}
\label{lem:ch7-cz-sharp-maximal}
若 $T$ 是 Calderón--Zygmund 算子, 则对每个 $s>1$,
\[
 M^\#(Tf)(x)\le C_sM(|f|^s)(x)^{1/s},
\]
其中 $M^\#$ 是 \eqref{eq:ch6-sharp-maximal-function} 定义的锐极大算子.
\end{Lemma}

\begin{Proof}
固定 $s>1$. 给定 $x$ 以及包含它的立方体 $Q$, 由命题~\ref{prop:ch6-bmo-equivalent-oscillation}, 只需找到常数 $a$, 使
\[
 \frac1{|Q|}\int_Q|Tf(y)-a|\,dy
 \le CM(|f|^s)(x)^{1/s}.
\]
如定理~\ref{thm:ch7-a1-characterization} 的证明那样, 写 $f=f_1+f_2$, 其中 $f_1=f\chi_{2Q}$. 令 $a=Tf_2(x)$, 则
\begin{equation}
\label{eq:ch7-sharp-maximal-split}
\begin{aligned}
 \frac1{|Q|}\int_Q|Tf(y)-a|\,dy
 &\le\frac1{|Q|}\int_Q|Tf_1(y)|\,dy\\
 &\quad+\frac1{|Q|}\int_Q|Tf_2(y)-Tf_2(x)|\,dy.
\end{aligned}
\end{equation}
由于 $s>1$, $T$ 在 $L^s(\mathbb R^n)$ 上有界. 因而
\[
\begin{aligned}
 \frac1{|Q|}\int_Q|Tf_1(y)|\,dy
 &\le\left(\frac1{|Q|}\int_Q|Tf_1(y)|^s\,dy\right)^{1/s}\\
 &\le C\left(\frac1{|Q|}\int_{2Q}|f(y)|^s\,dy\right)^{1/s}\\
 &\le2^{n/s}CM(|f|^s)(x)^{1/s}.
\end{aligned}
\]
用 \eqref{eq:ch5-standard-kernel-x-regularity} 估计 \eqref{eq:ch7-sharp-maximal-split} 右端第二项. 仍用 $\ell(Q)$ 表示 $Q$ 的边长, 得
\[
\begin{aligned}
 &\frac1{|Q|}\int_Q|Tf_2(y)-Tf_2(x)|\,dy\\
 &\quad\le\frac1{|Q|}\int_Q
 \left|\int_{\mathbb R^n\setminus2Q}[K(y,z)-K(x,z)]f(z)\,dz\right|dy\\
 &\quad\le\frac C{|Q|}\int_Q\int_{\mathbb R^n\setminus2Q}
 \frac{|y-x|^\delta}{|x-z|^{n+\delta}}|f(z)|\,dz\,dy\\
 &\quad\le C\ell(Q)^\delta\sum_{k=-1}^\infty
 \frac1{(2^k\ell(Q))^{n+\delta}}
 \int_{|x-z|<2^{k+1}\ell(Q)}|f(z)|\,dz\\
 &\quad\le CMf(x)\le CM(|f|^s)(x)^{1/s}.
\end{aligned}
\]
\end{Proof}

\hypertarget{chapter-07-page-156}{}
\begin{Lemma}
\label{lem:ch7-weighted-dyadic-sharp}
设 $w\in A_p$, $1\le p_0\le p<\infty$. 若 $M_df\in L^{p_0}(w)$, 则只要左端有限, 就有
\[
 \int_{\mathbb R^n}|M_df|^pw
 \le C\int_{\mathbb R^n}|M^\#f|^pw,
\]
其中 $M_d$ 是 \eqref{eq:ch2-dyadic-maximal} 定义的二进极大算子, $M^\#$ 是 \eqref{eq:ch6-sharp-maximal-function} 定义的锐极大算子.
\end{Lemma}

\begin{Proof}
这是引理~\ref{lem:ch6-dyadic-sharp-norm} 的带权版本, 证明几乎完全相同. 只需证明引理~\ref{lem:ch6-good-lambda} 的带权类似物: 对某个 $\delta>0$,
\[
\begin{aligned}
 &w(\{x\in\mathbb R^n:M_df(x)>2\lambda,\ M^\#f(x)\le\gamma\lambda\})\\
 &\qquad\le C\gamma^\delta
 w(\{x\in\mathbb R^n:M_df(x)>\lambda\}).
\end{aligned}
\]
集合 $\{x\in\mathbb R^n:M_df(x)>\lambda\}$ 可分解为两两不交的二进立方体, 因而只需对其中每个立方体 $Q$ 证明
\[
 w(\{x\in Q:M_df(x)>2\lambda,\ M^\#f(x)\le\gamma\lambda\})
 \le C\gamma^\delta w(Q).
\]
这立即由 \eqref{eq:ch6-good-lambda-cube} 和 $A_\infty$ 条件 \eqref{eq:ch7-ainfinity-condition} 得到: 前者是相应的 Lebesgue 测度不等式, 右端为 $C\gamma|Q|$.
\end{Proof}

\begin{Theorem}
\label{thm:ch7-cz-weighted-strong}
若 $T$ 是 Calderón--Zygmund 算子, 则对任意 $w\in A_p$, $1<p<\infty$, $T$ 在 $L^p(w)$ 上有界.
\end{Theorem}

\begin{Proof}
固定 $w\in A_p$. 由于具有紧支集的有界函数在 $L^p(w)$ 中稠密, 可设 $f$ 是这样的函数. 由推论~\ref{cor:ch7-ap-self-improvement}, 可取 $s>1$ 使 $w\in A_{p/s}$. 又因几乎处处有 $|Tf(x)|\le M_d(Tf)(x)$, 引理~\ref{lem:ch7-cz-sharp-maximal}, 引理~\ref{lem:ch7-weighted-dyadic-sharp} 和定理~\ref{thm:ch7-maximal-weighted-strong} 给出
\[
\begin{aligned}
 \int_{\mathbb R^n}|Tf|^pw
 &\le\int_{\mathbb R^n}M_d(Tf)^pw\\
 &\le C\int_{\mathbb R^n}M^\#(Tf)^pw\\
 &\le C\int_{\mathbb R^n}M(|f|^s)^{p/s}w\\
 &\le C\int_{\mathbb R^n}|f|^pw,
\end{aligned}
\]
其中暂时假设第二个积分有限. 为此只需证明 $Tf\in L^p(w)$. 若 $f$ 的支集包含于 $B(0,R)$, 则对 $\varepsilon>0$,
\[
\begin{aligned}
 \int_{|x|<2R}|Tf(x)|^pw(x)\,dx
 &\le\left(\int_{|x|<2R}w(x)^{1+\varepsilon}\,dx\right)^{1/(1+\varepsilon)}\\
 &\quad\times\left(\int_{|x|<2R}|Tf(x)|^{p(1+\varepsilon)/\varepsilon}\,dx\right)^{\varepsilon/(1+\varepsilon)}.
\end{aligned}
\]
由反向 Hölder 不等式可选择 $\varepsilon$, 使第一个积分有限; 第二个积分有限, 因为 $Tf\in L^q$, $1<q<\infty$.

\hypertarget{chapter-07-page-157}{}
另一方面, 若 $|x|>2R$, 则
\[
 |Tf(x)|=\left|\int_{\mathbb R^n}f(y)K(x,y)\,dy\right|
 \le C\int_{|y|<R}\frac{|f(y)|}{|x-y|^n}\,dy
 \le\frac{C\|f\|_\infty}{|x|^n}.
\]
因此
\[
\begin{aligned}
 \int_{|x|>2R}|Tf(x)|^pw(x)\,dx
 &\le C\sum_{k=1}^\infty
 \int_{2^kR<|x|<2^{k+1}R}\frac{w(x)}{|x|^{np}}\,dx\\
 &\le C\sum_{k=1}^\infty(2^kR)^{-np}w(B(0,2^{k+1}R)).
\end{aligned}
\]
由推论~\ref{cor:ch7-ap-self-improvement}, 存在 $q<p$ 使 $w\in A_q$. 由 \eqref{eq:ch7-weight-set-comparison} 对球的相应结论,
\[
 w(B(0,2^{k+1}R))\le C(n,R,w)2^{knq}.
\]
代入上式得到收敛级数. 因而 $Tf\in L^p(w)$, 证明完成.
\end{Proof}

\begin{Theorem}
\label{thm:ch7-cz-weighted-weak-one}
设 $T$ 是 Calderón--Zygmund 算子, $w\in A_1$. 则
\[
 w(\{x\in\mathbb R^n:|Tf(x)|>\lambda\})
 \le\frac C\lambda\int_{\mathbb R^n}|f|w.
\]
\end{Theorem}

\begin{Proof}
证明与定理~\ref{thm:ch5-calderon-zygmund} 的无权结论很相似, 并沿用其记号. 在高度 $\lambda$ 对 $f$ 作 Calderón--Zygmund 分解, 写 $f=g+b$. 只需分别估计 $w(\{|Tg|>\lambda\})$ 和 $w(\{|Tb|>\lambda\})$. 由命题~\ref{prop:ch7-ap-properties}, $w\in A_2$. 因而定理~\ref{thm:ch7-cz-weighted-strong} 给出
\[
\begin{aligned}
 w(\{|Tg|>\lambda\})
 &\le\frac1{\lambda^2}\int_{\mathbb R^n}|Tg|^2w
 \le\frac C{\lambda^2}\int_{\mathbb R^n}|g|^2w\\
 &\le\frac{2^nC}{\lambda}\int_{\mathbb R^n}|g|w.
\end{aligned}
\]

\hypertarget{chapter-07-page-158}{}
还需证明 $\int|g|w\le C\int|f|w$. 在 $\mathbb R^n\setminus\bigcup_jQ_j$ 上 $g=f$; 在每个 $Q_j$ 上, 由 $w\in A_1$,
\[
\begin{aligned}
 \int_{Q_j}|g|w
 &\le\int_{Q_j}\frac1{|Q_j|}\int_{Q_j}|f(y)|\,dy\,w(x)\,dx\\
 &=\int_{Q_j}|f(y)|\frac{w(Q_j)}{|Q_j|}\,dy
 \le C\int_{Q_j}|f(y)|w(y)\,dy.
\end{aligned}
\]
对坏函数部分, 由 \eqref{eq:ch7-weight-set-comparison},
\[
 w\left(\bigcup_jQ_j^*\right)
 \le\sum_jw(Q_j^*)\le C\sum_jw(Q_j)
 \le C\sum_j\frac{w(Q_j)}{|Q_j|}|Q_j|.
\]
而 $Q_j$ 的选取给出 $|Q_j|<\lambda^{-1}\int_{Q_j}|f|$, 所以上面的论证也控制了该和式.

令 $c_j$ 为 $Q_j$ 的中心. 因为 $b_j$ 在 $Q_j$ 上平均值为零,
\[
\begin{aligned}
 &w\left(\left\{x\in\mathbb R^n\setminus\bigcup_jQ_j^*:|Tb(x)|>\lambda\right\}\right)\\
 &\quad\le\frac C\lambda\sum_j\int_{\mathbb R^n\setminus Q_j^*}|Tb_j(x)|w(x)\,dx\\
 &\quad=\frac C\lambda\sum_j\int_{\mathbb R^n\setminus Q_j^*}
 \left|\int_{Q_j}[K(x,y)-K(x,c_j)]b_j(y)\,dy\right|w(x)\,dx.
\end{aligned}
\]
由 \eqref{eq:ch5-standard-kernel-y-regularity}, 最后一项不超过
\begin{equation}
\label{eq:ch7-weighted-bad-part}
 \frac C\lambda\sum_j\int_{Q_j}|b_j(y)|
 \left(\int_{\mathbb R^n\setminus Q_j^*}
 \frac{|y-c_j|^\delta}{|x-c_j|^{n+\delta}}w(x)\,dx\right)dy.
\end{equation}

\hypertarget{chapter-07-page-159}{}
如引理~\ref{lem:ch7-cz-sharp-maximal} 的证明所示, 括号内的量由 $CMw(y)$ 控制; 因 $w\in A_1$, 它又由 $Cw(y)$ 控制. 因而 \eqref{eq:ch7-weighted-bad-part} 不超过
\[
 \frac C\lambda\int_{\mathbb R^n}|b|w
 \le\frac C\lambda\int_{\mathbb R^n}(|f|+|g|)w
 \le\frac C\lambda\int_{\mathbb R^n}|f|w,
\]
其中最后一个不等式来自前面的论证.
\end{Proof}

最后回忆: 若 Calderón--Zygmund 算子 $T$ 满足
\begin{equation}
\label{eq:ch7-cz-truncation-limit}
 Tf(x)=\lim_{\varepsilon\to0}T_\varepsilon f(x),
\end{equation}
其中
\[
 T_\varepsilon f(x)=\int_{|x-y|>\varepsilon}K(x,y)f(y)\,dy,
\]
则称 $T$ 为 Calderón--Zygmund 奇异积分, 见第~\ref{sec:ch5-calderon-zygmund-singular-integrals}~节. 为判断极限 \eqref{eq:ch7-cz-truncation-limit} 对 $f\in L^p(w)$ 何时几乎处处存在, 需要考察相伴极大算子的加权范数不等式:
\[
 T^*f(x)=\sup_{\varepsilon>0}|T_\varepsilon f(x)|.
\]

\begin{Corollary}
\label{cor:ch7-maximal-cz-weighted}
若 $T$ 是 Calderón--Zygmund 算子, 则当 $1<p<\infty$ 且 $w\in A_p$ 时, $T^*$ 在 $L^p(w)$ 上有界; 当 $w\in A_1$ 时, $T^*$ 关于 $w$ 是弱型 $(1,1)$ 的.
\end{Corollary}

\begin{Proof}
这是定理~\ref{thm:ch5-maximal-cz-bounds} 的带权版本, 其证明依赖 Cotlar 不等式 \eqref{eq:ch5-cotlar-type-inequality}:
\[
 T^*f(x)\le C_\nu\left(M(|Tf|^\nu)(x)^{1/\nu}+Mf(x)\right),
 \qquad0<\nu<1.
\]
当 $1<p<\infty$ 时, 结论立即由定理~\ref{thm:ch7-maximal-weighted-strong} 和定理~\ref{thm:ch7-cz-weighted-strong} 得到. 当 $p=1$ 时, 按定理~\ref{thm:ch5-maximal-cz-bounds} 的证明进行, 并使用以下事实: 引理~\ref{lem:ch5-kolmogorov-weak-estimate} 在 Lebesgue 测度换成任意正 Borel 测度后仍成立; 极大函数的相应覆盖估计在 Lebesgue 测度换成 $w\,dx$ 且 $w\in A_1$ 后仍成立; 以及由定理~\ref{thm:ch7-maximal-weighted-weak}, $M$ 关于 $w$ 是弱型 $(1,1)$ 的.
\end{Proof}
